\chapter{Zdrojové kódy měřicích skriptů}
\label{sec:Appendix1}

Příloha obsahuje skripty, které vykonávají vlastní síťovou práci přístroje. Každý z nich je samostatným nástrojem spustitelným i z příkazové řádky ve jmenném prostoru zvoleného portu. Přehled všech souborů programového vybavení je uveden v Příloze \ref{sec:Appendix2}.

\section*{Příprava testovacích portů (setup\_network.sh)}
\lstinputlisting[language=bash, label=src:SetupNetwork, caption={Přesun portů do jmenných prostorů, nastavení adres a spouštění serveru dnsmasq}]{SourceCodes/setup_network.sh}

\section*{Analyzátor provozu (sniff\_stats.py)}
\lstinputlisting[language=python, label=src:SniffStats, caption={Zachytávání nástrojem tcpdump s filtrem BPF a výpočet statistik provozu}]{SourceCodes/sniff_stats.py}

\section*{Generátor paketů (pkt\_gen.py)}
\lstinputlisting[language=python, label=src:PktGen, caption={Sestavení rámců knihovnou Scapy a jejich odeslání zadanou rychlostí}]{SourceCodes/pkt_gen.py}

\section*{Měření propustnosti, ztrátovosti a rozptylu zpoždění (perf\_test.py)}
\lstinputlisting[language=python, label=src:PerfTest, caption={Řízení nástroje iperf3, rozmítání velikostí rámce a vyhodnocení výsledku}]{SourceCodes/perf_test.py}

\section*{Audit zpráv Router Advertisement (ra\_audit.py)}
\lstinputlisting[language=python, label=src:RaAudit, caption={Vyžádání, rozbor a posouzení zprávy Router Advertisement}]{SourceCodes/ra_audit.py}

\section*{Vyhledání sousedů protokolu IPv6 (neigh\_scan.py)}
\lstinputlisting[language=python, label=src:NeighScan, caption={Vyjmenování sousedů a odvození jejich adres postupem EUI-64}]{SourceCodes/neigh_scan.py}

\section*{Skenování protokolem ARP (arp\_scan.py)}
\lstinputlisting[language=python, label=src:ArpScan, caption={Vyhledání aktivních stanic v rozsahu adres portu}]{SourceCodes/arp_scan.py}
